\subsection{PROPERTIES OF INTERVALS IN R}

PROPOSITION 2. Every closed (resp. open) interval in R is a closed (resp. open) set in R.

The sets \([a, \rightarrow[ = a + \mathbf{R}_+ \text{ and } ]\leftarrow, a] = a - \mathbf{R}_+\) are obtained by translation from \(\mathbf{R}_+\) and \(-\mathbf{R}_+\) respectively and are therefore closed (Chapter III, § 1, no. 1); the sets \(]\leftarrow, a[\) and \(]a, \rightarrow[\) , which are their complements, are open; finally, the closed interval \([a, b]\) (resp. the open interval \(]a, b[\) ) is the intersection of \([a, \rightarrow[\) and \(]\leftarrow, b]\) (resp. of \(]a, \rightarrow[\) and \(]\leftarrow, b[\) ) and is therefore a closed (resp. open) set.

The closed intervals \([-a, +a]\) (a \textgreater{} 0) in R are therefore neighbourhoods of o. Let us show that they form a fundamental system of neighbourhoods of o as a runs through \(R_{+}^{*}\) . For this it is enough to 'establish the following proposition:

PROPOSITION 3. As r runs through the set of rational numbers \textgreater0, the intervals \(s_{r} = [-r, +r]\) in R form a fundamental system of neighbourhoods of 0.

By Proposition 7 of Chapter III, § 3, no. 4 we obtain a fundamental system of neighbourhoods of o in R by taking the closures in R of the intervals

\(S_{r} \cap Q = [-r, +r]\) of \(Q\) . The proof will be complete if we show that \(S_{r}\) is the closure of \(S_{r} \cap Q\) . Now \(S_{r}\) is closed in \(R\) , and we need therefore only prove that, if \(x\) is a real number such that \(-r < x < r\) , then \(x\) is in the closure of \(S_{r} \cap Q\) . The interval \(]-r, +r[\) is an open set in \(R\) and therefore for all sufficiently small neighbourhoods \(V\) of \(o\) in \(R\) we have \(x + V \subset ]-r, +r[\) ; but \(Q\) being dense in \(R\) , there is a rational number \(r' \in x + V\) , so that \(-r < r' < r\) and therefore \(r' \in S_{r} \cap Q\) .

COROLLARY. Every point of the real line has a countable fundamental system of neighbourhoods.

PROPOSITION 4. If \((x, y)\) is any pair of real numbers such that \(x < y\) , there is a rational number \(r\) such that \(x < r < y\) .

Since \(\mathbf{Q}\) is dense in \(\mathbf{R}\) , it is enough to show that \(]x, y[\) is not empty; by translation we may assume \(x = 0\) and \(y > 0\) . Now \(\mathbf{R}\) is a Hausdorff space and therefore, by Proposition 3, there is a rational number \(r > 0\) such that \(y \notin [-r, +r]\) , and this implies that \(0 < r < y\) .

PROPOSITION 5. Let I be any interval in R. Then the topology induced on I by the topology of R is generated by the open intervals of I (where I is considered as linearly ordered by the relation \(x \leqslant y\) ).

Every open interval of I is the trace on I of an open interval of R. This is clear for a bounded interval, and the unbounded interval \(]a, \rightarrow[\) of I is the trace of the unbounded interval \(]a, \rightarrow[\) of R. We may therefore restrict ourselves to the case I = R; but in this case the result follows from Proposition 3, since every neighbourhood of a point \(x \in R\) contains an open interval \(]x - a, x + a[\) .

Remark. If A is a dense subset of R, the topology of R is generated by the open intervals whose end-points belong to A. For if

\[
] x - a, \quad x + a [
\]

is an open interval containing \(x\) , there exist two points \(y, z\) of \(A\) such that \(x - a < y < x\) and \(x < z < x + a\) ; hence \(]y, z[\) contains \(x\) and is contained in \(]x - a, x + a[\) . This proof shows, moreover, that the intervals under consideration form a base (Chapter I, § 1, no. 3) of the topology of \(R\) . In particular, if we take \(A = Q\) , we see that the topology of \(R\) has a countable base.

